\documentclass[serif, 12pt]{beamer}

\usepackage{graphicx} % Allows including images
\usepackage{booktabs} % Allows the use of \toprule, \midrule and \bottomrule in tables

\usepackage{color}

%\usepackage{algorithm2e}
\usepackage{hyperref}
\usepackage{algorithm,algorithmic}
\usepackage{changepage}
\usepackage{pgfplots}
\usepackage{pgfplotstable}
\usepgfplotslibrary{statistics}

\newcommand*\mat[1]{ \begin{pmatrix} #1 \end{pmatrix}}
\newcommand*\arr[1]{ \begin{bmatrix} #1 \end{bmatrix}}
\newcommand*\V[1]{ \boldsymbol{#1}}

\newcommand*\D{\textcolor{violet}{D}}
\newcommand*\T{\textcolor{blue}{T}}

\setbeamertemplate{navigation symbols}{}%remove navigation symbols

\setbeamerfont{page number in head/foot}{size=\small}
\setbeamertemplate{footline}[frame number]

\title{Approximated PCA}
\subtitle{Experiment 8}

\author{Rodrigo Arias} % Your name
\date{\today} % Date, can be changed to a custom date

\begin{document}

\begin{frame}
	\titlepage
\end{frame}

%------------------------------------------------

\begin{frame}

\frametitle{Objective}

The objective of the experiment 8 is to analyze the behaviour of concatenating 
rounding operations.

\vspace{1em}
The rounding effects can be analyzed both theoretically and experimentally by 
using MPFR as the \textsl{real} result.

\vspace{1em}
Perform a set of experiments, and draw the error distribution as the number of 
operations increase.

\end{frame}

%------------------------------------------------

\begin{frame}

\frametitle{Notation}
\begin{center}
\begin{tabular}{l  l}
$x$ 											& Real number \\
$\tilde x = fl(x)$		 		& Approximation of $x$ \vspace{1em} \\
$\Delta = x-\tilde{x}$ 		& Rounding error, or simply error \\
$|\Delta|$								& Absolute error \vspace{1em} \\
$\delta = \Delta / |x|$		& Signed relative error \\
$|\delta| = |\Delta|/|x|$	& Relative error \vspace{1em} \\
$\epsilon = 2^{-b}$				& Maximum relative error. $|\delta| \le \epsilon$ \\
\end{tabular}
\end{center}

\end{frame}
%------------------------------------------------

\begin{frame}

\frametitle{Simple rounding, experiment 8a}

We generate a uniform random number $x \in [0.5, 1]$, and we apply a rounding 
operation to store an aproximation in 10 bits.

$$ \tilde{x} = fl(x) $$

Then, we measure the error $\Delta = x - \tilde x$. After some runs, we get the 
experimental distribution of $\Delta$.

\end{frame}

%------------------------------------------------

\begin{frame}

\frametitle{Distribution of $\Delta$}

Let $x$ be a number from a uniform random distribution in $[0.5,1]$. Let 
$\tilde{x} = fl(x)$, and with $b=10$, $\epsilon \approx 9.766 \cdot 10^{-4}$
%%
%\\\vspace{1em}
%%
%And so, the error $\Delta = x_1 - \tilde{x}_1 = r_1 - fl(r_1)$
%%
%\\\vspace{1em}
%%
%The error $\Delta$ is distributed as a uniform r.v. But as n increases starts to 
%shape as a normal distribution.
%
%
\\\vspace{1em}
%
\includegraphics[width=.9\textwidth]{img/px.pdf}
%
\\\vspace{1em}
%
The probability that $x$ falls in the interval $\tilde x \pm \epsilon$ is 
constant, as each interval has the same size. Inside the interval, the error 
$\Delta$ is equivalent to a uniform distribution in $[-\epsilon/2, 
+\epsilon/2]$.

\end{frame}

%------------------------------------------------

\begin{frame}
	\centering
	\includegraphics[width=\textwidth]{img/half-prob.png}
\end{frame}


%------------------------------------------------

\begin{frame}

\frametitle{Experiment 8c, multiple operations}

Let $(r_i)$ be a sequence of uniformly distributed random numbers in $[0, 1]$, 
and
%
\begin{equation}
\begin{array}{ll}
\begin{split}
	x_0 &= 0 \\
	x_i &= x_{i-1} + r_i
\end{split}
&\quad
\begin{split}
	\tilde{x}_0 &= fl(0) = 0 \\
	\tilde{x}_i &= fl(\tilde{x}_{i-1} + r_i)
\end{split}
\\
\end{array}
\end{equation}
%
Then we observe the error $\Delta = x_n - \tilde{x}_n$ as $n$ grows.

\end{frame}

%------------------------------------------------

\foreach \n in {1,2,3,4,5,6,7,8,9,10,20,30,40,50,60,70,80,90,100}
{
	\begin{frame}
%	\frametitle{Generated slide at $n=\n$}
	\centering
	\includegraphics[width=\textwidth]{img/8c-gsum-err-100000-\n-10.png}
	\end{frame}
}

%------------------------------------------------

\foreach \n in {1,2,3,4,5,6,7,8,9,10,20,30,40,50,60,70,80,90,100}
{
	\begin{frame}
%	\frametitle{Generated slide at $n=\n$}
	\centering
	\includegraphics[width=\textwidth]{img/8c-hist-err-100000-\n-10.png}
	\end{frame}
}

%------------------------------------------------

\foreach \n in {1,2,3,4,5,6,7,8,9,10,20,30,40,50,60,70,80,90,100}
{
	\begin{frame}
%	\frametitle{Generated slide at $n=\n$}
	\centering
	\includegraphics[width=\textwidth]{img/8c-qqnorm-err-100000-\n-10.png}
	\end{frame}
}

%------------------------------------------------

\begin{frame}

\frametitle{Distribution of $\Delta$}

As $n$ grows, the error $\Delta$ seems to approach a normal distribution.

\vspace{1em}
The mean $\mu$ is centered at 0, but the variance $\sigma^2$ seems to grow.

\vspace{1em}
After 1 operation, the expected theorical variance is
$$ \sigma^2 = \frac{1}{12} \epsilon^2 $$

\vspace{1em}
With 2 operations,
$$ \sigma^2 = \frac{1}{7} \epsilon^2 $$

\vspace{1em}
The relation between $\sigma^2$ and $n$ is not yet known.

\end{frame}

\end{document}
